% Generated draft: requires mathematical and source review.
\documentclass[11pt]{article}
\usepackage[margin=25mm]{geometry}
\usepackage{fontspec}
\setmainfont{Latin Modern Roman}
\usepackage{amsmath,amssymb,mathtools,unicode-math}
\setmathfont{Latin Modern Math}
\usepackage{enumitem,needspace,xurl,hyperref}
\setlength{\parindent}{0pt}
\setlength{\parskip}{5pt}
\setlength{\emergencystretch}{3em}
\newcommand{\N}{\mathbb{N}}
\newcommand{\Z}{\mathbb{Z}}
\newcommand{\Q}{\mathbb{Q}}
\newcommand{\R}{\mathbb{R}}
\newcommand{\C}{\mathbb{C}}
\newcommand{\F}{\mathbb{F}}
\newcommand{\A}{\mathbb{A}}
\newcommand{\PP}{\mathbb{P}}
\newcommand{\E}{\mathbb{E}}
\newcommand{\eps}{\varepsilon}
\newcommand{\dd}{\,\mathrm{d}}
\begin{document}
\section*{UnsolvedMath ID: 30001632}
\textbf{Source problem code:} catalogue entry
\subsection{Definitions and mathematical statement}
The supplied record asks two questions about the Takhtajan--Zograf Kähler form, but does not specify the underlying space, its topological parameters, or the analytic realization of the Laplacian. The following notation describes the requested objects conditionally; it does not supply these missing choices.

Let \(X\) denote the intended complex manifold, with complex structure \(J\), carrying the Takhtajan--Zograf Kähler form \(\omega_{\mathrm{TZ}}\). A Kähler form is a closed real differential form of type \((1,1)\) for which
\[
 g_{\mathrm{TZ}}(u,v)=\omega_{\mathrm{TZ}}(u,Jv)
\]
defines a positive definite Riemannian metric. This convention describes the associated metric; the supplied record does not give a construction or normalization of \(\omega_{\mathrm{TZ}}\).

\begin{enumerate}
\item Does \(\omega_{\mathrm{TZ}}\) admit a global representation formula on the intended space \(X\)? The source does not specify the permitted coordinates, constituent functions, or meaning of ``representation formula.'' Its supplied literature assessment mentions a global Fenchel--Nielsen-type formula, but does not establish that this is the exclusive intended target. Thus the required class of formulas remains unresolved.
\item Does the Laplace operator associated with \(g_{\mathrm{TZ}}\) admit a self-adjoint extension? To state this precisely, one must specify a Hilbert space \(H\) and a densely defined initial operator
\[
 \Delta_0\colon D_0\subset H\longrightarrow H.
\]
A self-adjoint extension means an operator \(A\colon D(A)\subset H\to H\) such that
\[
 D_0\subseteq D(A),\qquad
 Au=\Delta_0u\quad\text{for every }u\in D_0,
 \qquad A=A^*,
\]
where the last equality includes equality of operator domains. For a specified smooth Riemannian space, a natural Hilbert-space candidate is \(L^2(X,\mathrm{dvol}_{g_{\mathrm{TZ}}})\), but the record does not identify the intended space, initial domain, boundary conditions, treatment of singularities or degenerations, or sign convention. It also does not say whether a particular extension or uniqueness is intended. These choices must be recovered before the analytic question has a fully determined formulation.
\end{enumerate}

The supplied registry labels this problem \emph{partially solved}. Its dated assessment reports explicit potentials and representation formulas in important cases, including pointed rational moduli spaces, while reporting that no general global Fenchel--Nielsen-type formula has been established. This is supplied curation, not an independently verified mathematical or literature conclusion; it does not separately determine the status of the Laplacian question.
\subsection{Short English statement}
The problem asks for a global formula for the Takhtajan--Zograf Kähler form and a self-adjoint extension of its associated Laplacian. The underlying space, the meaning of the desired formula, and the initial operator domain are unspecified, so the draft retains both questions with these qualifications. The supplied partially solved label and special-case results require independent review.
\subsection{Sources}
\begin{itemize}
\item Supplied local UnsolvedMath registry record, permanent ID \(30001632\), source problem code \texttt{catalogue entry}, in the Oberwolfach Reports Open Problems collection. Title: \emph{Global Formulas and Laplacians for the Takhtajan--Zograf Form}. The record supplies proposed year \(2010\) and names no individual proposer.
\item Supplied report citation: \emph{Teichmüller Theory (2011)}, Oberwolfach Reports, pp.~3085--3157, DOI \texttt{10.4171/owr/2010/53}. Original OWR report: \url{https://doi.org/10.4171/owr/2010/53}. The supplied background describes Oberwolfach Reports as documenting mathematical workshops and their open-problem discussions.
\item Supplied literature source, labelled ``potentials for Takhtajan--Zograf metrics'': \url{https://www.sciencedirect.com/science/article/am/pii/S0001870816301670}. No author, article title, or theorem locator is supplied for this link.
\item Literature assessment date: \texttt{2026-08-21}. This is supplied dated curation derived from the OWR curation, not independently verified. It reports partial results, including explicit potentials and representation formulas for pointed rational moduli spaces, and the absence of an established general global Fenchel--Nielsen-type formula. The record explicitly requests independent verification of time-sensitive status and cited claims.
\item The registry supplies statement status \texttt{corrected\_verified} and says its mathematical content was checked against the official report, with notation, context, and question form normalized without changing the source-level claim. This verification assertion is part of the supplied provenance; no source was independently visited for this draft.
\end{itemize}
\end{document}
